\documentclass[11pt]{article}

\usepackage[margin=1in]{geometry}
\usepackage{amsmath,amssymb,amsthm}
\usepackage{booktabs}
\usepackage{microtype}
\usepackage{natbib}
\usepackage{xcolor}
\setlength{\emergencystretch}{3em}
\usepackage[colorlinks=true,linkcolor=blue!50!black,citecolor=blue!50!black,urlcolor=blue!50!black]{hyperref}

\newtheorem{proposition}{Proposition}
\newtheorem{corollary}[proposition]{Corollary}
\theoremstyle{definition}
\newtheorem{definition}{Definition}
\theoremstyle{remark}
\newtheorem{remark}{Remark}

\newcommand{\KL}{D_{\mathrm{KL}}}
\newcommand{\softmax}{\operatorname{softmax}}
\newcommand{\Tr}{\operatorname{Tr}}
\newcommand{\diag}{\operatorname{diag}}
\newcommand{\E}{\mathbb{E}}
\newcommand{\N}{\mathcal{N}}
\newcommand{\F}{\mathcal{F}}

\title{Template Routing:\\ The Block MLP as Inference over Learned Gaussian Concepts}
\author{Robert C. Dennis}
\date{Explanatory note, September 2026}

\begin{document}
\maketitle

\begin{abstract}
In the gauge-covariant Gaussian transformer, attention is derived as variational inference among token agents, while the residual multilayer perceptron (MLP) that follows each attention step is imposed without derivation. This note develops template routing, a replacement for that MLP that follows from the same variational principle. Each token agent is coupled, in addition to its peers in the sequence, to a bank of learned Gaussian concept agents that live outside every sequence and whose beliefs are set by training rather than by inference. A mixture prior over these concepts yields responsibilities proportional to $\pi_a\exp(-\KL(q_i\|p_a))$, a free-energy bound whose optimum is a soft minimum of divergences, and a mean update that reduces, for templates with a shared covariance, to a width-$M$ MLP whose pointwise activation is replaced by a softmax over units. We derive the responsibilities, the resulting mean updates with tied and untied values, the dependence on the agent's uncertainty that an ordinary MLP lacks, and a gauge-covariant form in reference coordinates. A worked example follows two tokens of one sentence through the routing. A structured posterior, exact when the attention output is read as a Gaussian message, turns each token's belief into a Gaussian mixture over template readings and yields a context-dependent ambiguity covariance. A social reading identifies the templates with a shared repertoire of norms and scripts: uncertain agents conform more, norms computed by the agents themselves settle at the weighted average of their adherents, and the templates are not meta-agents formed from particular coalitions. We then describe the implementation in \texttt{vfe3} and state what the construction does not establish.
\end{abstract}

\section{The problem this note addresses}

The evaluated architecture interleaves two kinds of step. The attention step is a consequence of the model: each token $i$ carries a Gaussian belief $q_i=\N(\mu_i,\Sigma_i)$, compares itself with its peers through a divergence, and receives the message
\begin{equation}
m_i=\sum_j\beta_{ij}L_{ij}\mu_j,
\qquad
\beta_{ij}\propto\pi_{ij}\exp\bigl(-D_{ij}/\tau\bigr),
\label{eq:attention}
\end{equation}
where $L_{ij}$ is the content link and $D_{ij}$ the divergence between query and key observables. The Gibbs form of $\beta_{ij}$ is derived from a free energy, and its covariance under frame changes follows from the transformation laws of the Gaussian data. The step that follows it is not derived. After normalization, each mean is updated by a coordinate MLP,
\begin{equation}
\mu_i\mapsto\mu_i+W_2\operatorname{GELU}(W_1\mu_i+b_1)+b_2,
\label{eq:mlp}
\end{equation}
with the Gaussian error linear unit of \citet{hendrycks2016}. The MLP carries a large share of the model's parameters, and the model has no account of why it should be there or what function it computes.

Two tempting explanations do not work. The first is that a group-valued link $\Theta_{ij}=\exp(A_{ij})$ in place of the matrix link $L_{ij}$ would supply the missing nonlinearity. It does not, because the message $\sum_j\beta_{ij}\Theta_{ij}\mu_j$ remains linear in the states $\mu_j$ for any link that does not itself depend on the states; the exponential acts on the Lie-algebra parameter, not on the belief. The second is that the MLP is a symmetry-breaking afterthought that could be removed. Removing it deletes most of the block's capacity, since a $K\times K$ operator acting on a $K$-vector cannot reproduce a map that expands to width $M\approx 4K$ to $12K$ and gates each unit.

Template routing takes a different route. The exponential that the model does contain lives in the routing weights of Eq.~\eqref{eq:attention}, and the MLP can be read as routing to a fixed set of learned targets. This reading is known in the transformer literature, where feed-forward layers behave as key--value memories \citep{geva2021} and can be replaced by attention over persistent learned vectors \citep{sukhbaatar2019}, and it connects attention to associative retrieval in modern Hopfield networks \citep{ramsauer2021}. What the present framework adds is a generative account of the targets, as Gaussian concept agents, and a derivation of the routing rule, the update, and its dependence on uncertainty from a free energy.

\section{Agents on two timescales}

The model already treats tokens as agents. Agent $i$ holds a belief $q_i=\N(\mu_i,\Sigma_i)$ on its fiber $E_i\cong\mathbb{R}^K$, with a frame $U_i$ relating local coordinates to a common reference space. Within a forward pass, agents exchange evidence through attention; the frame, the token priors, and the communication maps are learned across the corpus. Two sources of belief are therefore available to an agent during inference: its token-type prior, fixed once at the embedding, and its peers in the current sequence. The peers know only the present context. Nothing in the architecture lets an agent consult, in the middle of the computation, what the corpus has taught about states like its own.

\begin{definition}[Concept agents]
A bank of $M$ concept agents is a collection $\{(\pi_a,p_a,w_a)\}_{a=1}^M$, where $\pi_a>0$ with $\sum_a\pi_a=1$ is a prior weight, $p_a=\N(\nu_a,\Lambda_a)$ is a Gaussian template with $\Lambda_a\succ0$, and $w_a\in\mathbb{R}^K$ is a value. Concept agents are not tokens, are shared across all sequences, and hold beliefs that are parameters of the network, updated by training and fixed during a forward pass.
\end{definition}

The distinction between token agents and concept agents is one of timescale. Token beliefs are inferred on the fast timescale of a forward pass; concept beliefs are learned on the slow timescale of training, in the manner of a higher level of a hierarchical generative model whose states change slowly relative to the level below \citep{friston2008hierarchical}. The attention step answers the question \emph{what does this context say about me}. Template routing answers a second question, \emph{which known kind of thing am I, and what does that imply}. The rest of this note makes the second question precise; Section~\ref{sec:example} follows one sentence through the construction, and Section~\ref{sec:social} examines how far the reading of templates as shared social norms holds.

\section{The mixture prior and its free energy}

Let agent $i$ place a mixture prior on its latent state $k\in\mathbb{R}^K$,
\begin{equation}
p(k)=\sum_{a=1}^M\pi_a p_a(k),
\label{eq:mixture_prior}
\end{equation}
realized by a categorical latent $z\in\{1,\dots,M\}$ with $p(z=a)=\pi_a$ and $p(k\mid z=a)=p_a(k)$. The agent's belief over the joint latent $(k,z)$ is taken in the mean-field form $q_i(k)\rho_i(z)$, with $q_i$ given (it is the belief delivered by the preceding attention step) and $\rho_i$ a distribution on the $M$ labels. The variational free energy of the prior term is
\begin{equation}
\F_i[\rho_i]
=\E_{q_i\rho_i}\bigl[\log q_i(k)+\log\rho_i(z)-\log\pi_z-\log p_z(k)\bigr]
=\sum_a\rho_{ia}\Bigl[\KL(q_i\|p_a)+\log\frac{\rho_{ia}}{\pi_a}\Bigr].
\label{eq:free_energy}
\end{equation}

\begin{proposition}[Responsibilities]
\label{prop:responsibilities}
The minimizer of Eq.~\eqref{eq:free_energy} over the probability simplex is
\begin{equation}
\rho_{ia}^{*}=\frac{\pi_a\exp\bigl(-\KL(q_i\|p_a)\bigr)}{\sum_b\pi_b\exp\bigl(-\KL(q_i\|p_b)\bigr)},
\label{eq:responsibilities}
\end{equation}
and the minimum value is the soft minimum
\begin{equation}
\F_i^{*}=-\log\sum_a\pi_a\exp\bigl(-\KL(q_i\|p_a)\bigr).
\label{eq:softmin}
\end{equation}
Moreover, $\KL\bigl(q_i\|\sum_a\pi_ap_a\bigr)\le\F_i[\rho_i]$ for every $\rho_i$, so $\F_i^{*}$ is the tightest bound of this form on the divergence from the mixture prior.
\end{proposition}

\begin{proof}
Write $c_a=\KL(q_i\|p_a)-\log\pi_a$. Then $\F_i[\rho_i]=\sum_a\rho_{ia}(c_a+\log\rho_{ia})$. With $Z=\sum_be^{-c_b}$ and $\rho^{*}_{ia}=e^{-c_a}/Z$, this equals $\KL(\rho_i\|\rho_i^{*})-\log Z$, which is minimized exactly at $\rho_i=\rho_i^{*}$ with value $-\log Z$; that is Eqs.~\eqref{eq:responsibilities} and~\eqref{eq:softmin}. For the bound, Jensen's inequality applied to the concave logarithm gives, pointwise in $k$,
\begin{equation*}
-\log\sum_a\pi_ap_a(k)=-\log\sum_a\rho_{ia}\frac{\pi_ap_a(k)}{\rho_{ia}}
\le-\sum_a\rho_{ia}\log\frac{\pi_ap_a(k)}{\rho_{ia}}.
\end{equation*}
Taking $\E_{q_i}$ and adding $\E_{q_i}[\log q_i]$ to both sides yields $\KL(q_i\|\sum_a\pi_ap_a)\le\F_i[\rho_i]$.
\end{proof}

The entropy of $q_i$ never appears in Eq.~\eqref{eq:responsibilities}: in the cross-entropy $\E_{q_i}[-\log p_a]=\KL(q_i\|p_a)+H(q_i)$ it is the same for every template and cancels in the normalization. The routing rule is therefore the same Gibbs form as attention, with divergence to a concept in place of divergence to a peer and unit temperature fixed by the derivation rather than chosen. For a single-component belief, Eq.~\eqref{eq:softmin} coincides with the variational approximation to the divergence between a Gaussian and a Gaussian mixture introduced by \citet{hershey2007}. The bound is tight when the posterior over labels, $p(z\mid k)$, is constant over the support of $q_i$, which is the case when one template dominates or when $q_i$ is concentrated.

\section{Explicit logits for diagonal templates}

The implementation uses diagonal templates, $\Lambda_a=\diag(\lambda_{a1},\dots,\lambda_{aK})$, which keeps the parameter count linear in $K$ per template. For a full belief covariance $\Sigma_i$, the Gaussian divergence is
\begin{equation}
\KL(q_i\|p_a)=\frac12\Bigl[\Tr(\Lambda_a^{-1}\Sigma_i)+(\mu_i-\nu_a)^{\top}\Lambda_a^{-1}(\mu_i-\nu_a)-K+\log\det\Lambda_a-\log\det\Sigma_i\Bigr].
\label{eq:gaussian_kl}
\end{equation}
The terms $-K$ and $-\log\det\Sigma_i$ do not depend on $a$ and cancel in Eq.~\eqref{eq:responsibilities}. With a diagonal $\Lambda_a$, the trace reads only the diagonal of $\Sigma_i$, so the logits are
\begin{equation}
\ell_{ia}=\log\pi_a-\frac12\sum_{k=1}^K\Bigl[\frac{(\Sigma_i)_{kk}+(\mu_{ik}-\nu_{ak})^2}{\lambda_{ak}}+\log\lambda_{ak}\Bigr],
\qquad
\rho^{*}_{ia}=\softmax_a(\ell_{ia}).
\label{eq:logits}
\end{equation}
Two consequences follow directly. The logits are exact up to an $a$-independent constant, not an approximation to the divergence, and they require no determinant or inverse of the belief covariance. And a full covariance and its diagonal produce identical responsibilities against diagonal templates, so the construction is indifferent to the storage format of $\Sigma_i$. Expanding the square writes the quadratic part as three matrix products,
\begin{equation}
\sum_k\frac{(\Sigma_i)_{kk}+\mu_{ik}^2}{\lambda_{ak}}
-2\sum_k\frac{\mu_{ik}\nu_{ak}}{\lambda_{ak}}
+\sum_k\frac{\nu_{ak}^2}{\lambda_{ak}},
\label{eq:matmul_form}
\end{equation}
each a contraction of a $K$-vector against an $M\times K$ table, which is how the logits are computed.

\section{Relation to the MLP}

Consider templates that share one covariance, $\Lambda_a=\Lambda$ for every $a$, with an arbitrary belief covariance $\Sigma_i$. Then $\Tr(\Lambda^{-1}\Sigma_i)$, $\log\det\Lambda$, and $\mu_i^{\top}\Lambda^{-1}\mu_i$ are all independent of $a$ and cancel in the softmax, leaving
\begin{equation}
\ell_{ia}=\nu_a^{\top}\Lambda^{-1}\mu_i-\frac12\nu_a^{\top}\Lambda^{-1}\nu_a+\log\pi_a+\text{const}_i.
\label{eq:affine_logits}
\end{equation}
These are the pre-activations of a linear layer: the rows of the weight matrix are the whitened template means $\Lambda^{-1}\nu_a$, and the bias of each unit is fixed by its template's norm and prior weight. If the update is $\mu_i\mapsto\mu_i+\sum_a\rho^{*}_{ia}w_a$ (Section~\ref{sec:update}), the values $w_a$ are the columns of the output matrix. Template routing in this regime is therefore a width-$M$ MLP whose pointwise activation has been replaced by a softmax across units. Table~\ref{tab:correspondence} records the correspondence.

\begin{table}[h]
\centering
\begin{tabular}{@{}lll@{}}
\toprule
 & Coordinate MLP, Eq.~\eqref{eq:mlp} & Template routing \\
\midrule
Unit $a$ reads & row $a$ of $W_1$ & template $(\nu_a,\Lambda_a)$ \\
Similarity & $(W_1\mu_i+b_1)_a$ & $-\KL(q_i\|p_a)+\log\pi_a$ \\
Gate & $\operatorname{GELU}$, independent per unit & $\softmax$ across units \\
Unit $a$ writes & column $a$ of $W_2$ & value $w_a$ \\
Depends on $\Sigma_i$ & no & yes, through $\diag\Sigma_i$ \\
Parameters per layer & $2MK+M+K$ & $3MK+M$ \\
\bottomrule
\end{tabular}
\caption{Correspondence between the coordinate MLP and template routing at width $M$. The two coincide in form, up to the gate, when the templates share a covariance.}
\label{tab:correspondence}
\end{table}

The gate is the remaining difference, and it is smaller than it appears. The sigmoid-weighted linear unit satisfies the identity
\begin{equation}
x\sigma(x)=x\cdot\frac{e^{x}}{e^{x}+e^{0}},
\label{eq:silu}
\end{equation}
which is softmax routing over two options with logits $(x,0)$ and values $(x,0)$ \citep{elfwing2018,ramachandran2017}; GELU is close to $x\sigma(1.702x)$ \citep{hendrycks2016}. A GELU unit is thus a two-way competition between ``pass the pre-activation'' and ``pass nothing,'' decided by an exponential weight. Template routing replaces $M$ independent two-way competitions by a single $M$-way competition among concepts.

\begin{remark}[Independent gates]
\label{rem:independent}
Independent per-unit gates also have a generative reading. Give each concept its own binary latent that selects between the template $p_a$ with probability $\pi_a$ and a shared background density $p_0$ otherwise. The same argument as Proposition~\ref{prop:responsibilities}, applied to each two-component mixture separately, gives the gate
\begin{equation}
g_{ia}=\sigma\Bigl(\log\frac{\pi_a}{1-\pi_a}-\KL(q_i\|p_a)+\KL(q_i\|p_0)\Bigr),
\end{equation}
which is the two-option competition of Eq.~\eqref{eq:silu} with log-odds given by a divergence difference, and the block output is $\sum_ag_{ia}w_a$. The choice between the softmax of Eq.~\eqref{eq:responsibilities} and these gates is a modeling statement: one concept explains the agent, or several concepts may explain it at once. The implementation uses the softmax. The independent-gate variant is not implemented.
\end{remark}

\section{From responsibilities to a mean update}
\label{sec:update}

The responsibilities say which concepts the agent recognizes itself as; a mean update must say what that recognition does to the belief. Two constructions are available, and they differ in what the values $w_a$ are allowed to be.

\paragraph{Tied values: descent on the concept free energy.}
Treat the belief mean as free and descend the optimal free energy of Eq.~\eqref{eq:softmin}. By the envelope theorem, the gradient of the soft minimum is the responsibility-weighted gradient of the divergences,
\begin{equation}
\frac{\partial\F_i^{*}}{\partial\mu_i}=\sum_a\rho^{*}_{ia}\Lambda_a^{-1}(\mu_i-\nu_a).
\end{equation}
The Fisher metric of the Gaussian location family is $\Sigma_i^{-1}$, so the natural-gradient step with rate $\eta$ is
\begin{equation}
\mu_i\mapsto\mu_i-\eta\Sigma_i\sum_a\rho^{*}_{ia}\Lambda_a^{-1}(\mu_i-\nu_a).
\label{eq:tied_update}
\end{equation}
Each agent is pulled toward the concepts it recognizes, weighted by how strongly it recognizes each one. With $\Lambda_a=\lambda I$ and a Euclidean step of size $\lambda$, the update is $\mu_i\mapsto\sum_a\rho^{*}_{ia}\nu_a$ with $\rho^{*}_{ia}=\softmax_a(\nu_a^{\top}\mu_i/\lambda-\|\nu_a\|^2/2\lambda+\log\pi_a)$, the retrieval step of a modern Hopfield network with inverse temperature $1/\lambda$ and stored patterns $\nu_a$ \citep{ramsauer2021}; the norm term vanishes for patterns of equal norm. In this construction the value written by concept $a$ is tied to its key: recognition moves the belief toward the prototype. The operation cleans up or denoises a state, and it cannot write information that the prototype does not already contain.

\paragraph{Untied values: concepts that predict.}
An MLP does more than clean up a state, because its output columns are independent of its input rows: a recognized feature can add a direction unrelated to the feature itself. For template routing to do the same, a concept must predict something beyond its own location. Let each concept be a joint Gaussian over the agent's state $k$ and an increment $y\in\mathbb{R}^K$,
\begin{equation}
p_a(k,y)=\N\left(\begin{pmatrix}k\\ y\end{pmatrix};\begin{pmatrix}\nu_a\\ w_a\end{pmatrix},\begin{pmatrix}\Lambda_a&\Lambda_aB_a^{\top}\\ B_a\Lambda_a&\Gamma_a\end{pmatrix}\right),
\end{equation}
with a regression block $B_a$ and $\Gamma_a-B_a\Lambda_aB_a^{\top}\succ0$, so that the joint covariance is positive definite. For a point state $k$, conditioning the mixture $\sum_a\pi_ap_a(k,y)$ on $k$ gives the standard Gaussian mixture regression \citep{ghahramani1994,jacobs1991},
\begin{equation}
\E[y\mid k]=\sum_ar_a(k)\bigl(w_a+B_a(k-\nu_a)\bigr),
\qquad
r_a(k)=\frac{\pi_a\N(k;\nu_a,\Lambda_a)}{\sum_b\pi_b\N(k;\nu_b,\Lambda_b)}.
\label{eq:gmr}
\end{equation}
Each concept is a local linear expert; the responsibilities decide which experts speak. Setting $B_a=0$ leaves the piecewise-constant readout $\sum_ar_a(k)w_a$, and the implemented block is
\begin{equation}
\mu_i\mapsto\mu_i+\sum_a\rho^{*}_{ia}w_a,
\qquad
\Sigma_i\mapsto\Sigma_i.
\label{eq:implemented}
\end{equation}

Equation~\eqref{eq:implemented} uses the variational responsibilities $\rho^{*}_{ia}$ of the belief, not the point responsibilities $r_a(k)$. The two agree in the limit of a concentrated belief: as $\Sigma_i\to0$, the $a$-dependent part of $\KL(q_i\|p_a)$ tends to $-\log\N(\mu_i;\nu_a,\Lambda_a)$, so $\rho^{*}_{ia}\to r_a(\mu_i)$. For a belief of finite width, $\sum_a\rho^{*}_{ia}w_a$ is not the exact expectation $\E_{q_i}\E[y\mid k]=\sum_a\E_{q_i}[r_a(k)]w_a$; it is the readout of the optimal mean-field label posterior of Proposition~\ref{prop:responsibilities}. The implemented update is therefore a declared variational construction, exact at the point-belief limit, and not an exact conditional expectation for general $\Sigma_i$. The regression blocks $B_a$ are a direct extension, strictly richer than a GELU MLP of the same width, and are not implemented.

\section{What the covariance does}

Equation~\eqref{eq:logits} contains the belief covariance through $\sum_k(\Sigma_i)_{kk}/\lambda_{ak}$, a term that has no counterpart in Eq.~\eqref{eq:mlp}. Its effect is to route uncertain agents toward broad concepts. For an isotropic template $\Lambda_a=\lambda I$, an isotropic belief $\Sigma_i=sI$, and a fixed squared distance $d^2=\|\mu_i-\nu_a\|^2$, the divergence is $\frac12[(Ks+d^2)/\lambda+K\log\lambda]$ up to $a$-independent terms, and it is minimized over the template width at
\begin{equation}
\lambda^{*}=s+\frac{d^2}{K}.
\end{equation}
An agent that is confident ($s$ small) prefers a template that is both near and sharp; an agent that is uncertain ($s$ large) pays a large trace penalty $Ks/2\lambda$ against any sharp template and prefers broad ones whose width matches its own. The same token mean is therefore routed differently depending on how certain the agent is, which is behavior a coordinate MLP cannot express because it never reads $\Sigma_i$. The template widths also play the role of a temperature: scaling every $\Lambda_a$ by $c$ scales the $a$-dependent quadratic terms by $1/c$. This is why no separate temperature parameter is introduced, and why the unit temperature of Proposition~\ref{prop:responsibilities} is not a restriction.

In the implemented block the covariance is read but not written: Eq.~\eqref{eq:implemented} passes $\Sigma_i$ through unchanged, consistent with the evaluated architecture's treatment of covariance as a routing state. A covariance update from the concept step, for instance by moment matching of the mixture readout, is a possible extension.

\section{Gauge covariance}

The frame field acts on agent data by $U_i\mapsto g_iU_i$, $\mu_i\mapsto g_i\mu_i$, and $\Sigma_i\mapsto g_i\Sigma_ig_i^{\top}$ with $g_i\in\mathrm{GL}(K)$. The reference coordinates
\begin{equation}
\xi_i=U_i^{-1}\mu_i,
\qquad
C_i=U_i^{-1}\Sigma_iU_i^{-\top}
\end{equation}
are invariant under this action. Place the concepts in the reference space and let agent $i$ see concept $a$ through its own frame, as the pushforward $(U_i)_*p_a=\N(U_i\nu_a,U_i\Lambda_aU_i^{\top})$.

\begin{proposition}[Covariant template routing]
\label{prop:covariant}
With responsibilities computed from $\KL\bigl(q_i\|(U_i)_*p_a\bigr)$ and message $m_i=U_i\sum_a\rho^{*}_{ia}w_a$, the responsibilities are invariant and the message transforms as $m_i\mapsto g_im_i$, so the update $\mu_i\mapsto\mu_i+m_i$ has the same covariance law as the attention residual.
\end{proposition}

\begin{proof}
The divergence is invariant under a common invertible affine change of variables, so $\KL\bigl(q_i\|(U_i)_*p_a\bigr)=\KL\bigl(\N(\xi_i,C_i)\|p_a\bigr)$, which depends only on invariant data. The responsibilities are functions of these divergences and the fixed $\pi_a$, hence invariant. The message is $U_i$ applied to an invariant vector and so transforms as $g_iU_i(\cdot)=g_im_i$.
\end{proof}

In this form the concept bank is a shared, frame-free library that every agent reads in its own coordinates, which is the same pattern by which attention compares agents through their frames. Diagonal templates are well defined here because the reference space has fixed coordinates. The evaluated architecture applies its normalization and MLP in local coordinates, and the implemented template block follows it: it computes Eq.~\eqref{eq:logits} from $(\mu_i,\Sigma_i)$ directly, which is Proposition~\ref{prop:covariant} with every $U_i$ replaced by the identity. It is therefore no more and no less covariant than the MLP it replaces. The covariant variant would reuse the realized canonical frame already available to the canonical-frame MLP mode and is not implemented.

\section{Relation to attention}

Template routing is attention in which the keys are concept templates and the values are concept outputs. Setting the two messages side by side,
\begin{equation}
m_i^{\mathrm{peers}}=\sum_j\beta_{ij}L_{ij}\mu_j,
\qquad
m_i^{\mathrm{concepts}}=\sum_a\rho^{*}_{ia}w_a,
\end{equation}
the peer message transports what another agent believes, while the concept message writes what the corpus associates with a kind of state. Both weights are Gibbs distributions over divergences. The implementation keeps the two normalizations separate and applies the concept step after the attention step and normalization, in the position of the MLP. A single softmax over the union of peers and concepts is the alternative studied by \citet{sukhbaatar2019}, who appended persistent vectors to every attention layer and removed the feed-forward sublayer; it corresponds to one mixture prior over both sources of evidence rather than two sequential inferences.

\section{A worked example: one sentence, token by token}
\label{sec:example}

This section follows two tokens of the sentence ``a man went to the bank to deposit \dots'' through one template step. The numbers are a hand-built illustration in $K=2$ with $M=3$, chosen so that every quantity can be checked by hand; a trained layer has $K$ in the tens, hundreds of templates, and templates without human labels. Read the first coordinate as ``finance'' and the second as ``terrain and water.'' Table~\ref{tab:toy_templates} lists the templates, all with uniform prior weight $\pi_a=1/3$.

\begin{table}[h]
\centering
\begin{tabular}{@{}llll@{}}
\toprule
Template & Mean $\nu_a$ & Widths $\lambda_a$ & What its value $w_a$ adds \\
\midrule
$a_1$ ``money institution'' & $(2,0)$ & $(0.5,0.5)$, sharp & toward \emph{money, check, cash, his} \\
$a_2$ ``riverbank'' & $(0,2)$ & $(0.5,0.5)$, sharp & toward \emph{river, water, sat, grass} \\
$a_3$ ``generic place noun'' & $(0.5,0.5)$ & $(3,3)$, broad & small generic push toward \emph{to, and, .} \\
\bottomrule
\end{tabular}
\caption{Toy template bank for the worked example. The labels are for the reader; the model sees only the parameters.}
\label{tab:toy_templates}
\end{table}

\paragraph{The token ``bank.''}
Attention is causal, so at position six ``bank'' has seen only ``a man went to the.'' Nothing in that prefix separates the financial reading from the riverside one, so after attention its belief lies between the two senses and is uncertain: $\mu=(1,1)$ and $\Sigma=2I$. Equation~\eqref{eq:logits} gives the logits and responsibilities in the first row of Table~\ref{tab:toy_routing}. The sharp templates are each at squared distance $2$ from the belief, and the trace term $\sum_k\Sigma_{kk}/\lambda_{ak}=8$ penalizes them heavily; the broad template pays a trace term of only $4/3$. The uncertain token routes almost entirely to the generic template and receives a generic update, which is the correct behavior when the evidence for either sense is absent. The representation of ``bank'' at its own position never sees ``deposit''; it exists to predict the seventh token, and disambiguation happens at later positions.

\paragraph{The token ``deposit.''}
At position eight, attention lets ``deposit'' read its past, including ``bank.'' The combination is unambiguous, so the belief is displaced toward finance and is confident: $\mu=(1.8,0.2)$ and $\Sigma=0.3I$. The second row of Table~\ref{tab:toy_routing} shows the result. The token recognizes itself as an instance of $a_1$ with responsibility $0.82$, and the update $\mu\mapsto\mu+0.82w_1+0.18w_3$ moves its belief toward the region the decoder reads as \emph{money}, \emph{his}, or \emph{the check}, which is what should follow ``deposit.''

\paragraph{The same mean at lower confidence.}
The third row keeps $\mu=(1.8,0.2)$ but sets $\Sigma=2I$, as it would be after a muddled context. The routing reverses toward the hedging template, $0.21$ to $a_1$ and $0.79$ to $a_3$. A coordinate MLP reads only $\mu$ and would return the same output in both rows. Template routing distinguishes ``in financial territory and sure of it'' from ``in financial territory and guessing,'' which is the behavior derived in the preceding sections.

\begin{table}[h]
\centering
\begin{tabular}{@{}lcccccc@{}}
\toprule
Token state & $\ell_{1}$ & $\ell_{2}$ & $\ell_{3}$ & $\rho^{*}_{1}$ & $\rho^{*}_{2}$ & $\rho^{*}_{3}$ \\
\midrule
``bank,'' $\mu=(1,1)$, $\Sigma=2I$ & $-5.31$ & $-5.31$ & $-1.85$ & $0.030$ & $0.030$ & $0.941$ \\
``deposit,'' $\mu=(1.8,0.2)$, $\Sigma=0.3I$ & $0.01$ & $-6.39$ & $-1.50$ & $0.818$ & $0.001$ & $0.181$ \\
``deposit,'' $\mu=(1.8,0.2)$, $\Sigma=2I$ & $-3.39$ & $-9.79$ & $-2.06$ & $0.210$ & $0.000$ & $0.790$ \\
\bottomrule
\end{tabular}
\caption{Logits of Eq.~\eqref{eq:logits} (dropping the common $\log\pi_a$) and responsibilities of Eq.~\eqref{eq:responsibilities} for the toy bank of Table~\ref{tab:toy_templates}, computed exactly and rounded.}
\label{tab:toy_routing}
\end{table}

The example shows the division of labor within one layer. Attention supplies context: it is how ``deposit'' learns that a bank was mentioned. The template step supplies corpus knowledge: the fact that money words follow ``deposit'' near finance is stored in no token of this sentence, only in $(\nu_1,\lambda_1,w_1)$. Neither step alone produces the right prediction. Training shapes the bank. Whenever $a_1$ fires and the next token is ``his'' or ``money,'' the gradient of the next-token loss turns $w_1$ toward that prediction, draws $\nu_1$ toward the beliefs that fire it, adjusts $\lambda_1$ to their spread, and raises $\pi_1$ if the situation is common.

\section{A structured posterior: multimodal beliefs from the template step}
\label{sec:structured}

The construction so far keeps one Gaussian per token and puts all multimodality in the categorical label: the posterior $q_i(k)\rho_i(z)$ is mean-field, so the belief about the state does not depend on which template is responsible. This section develops the structured alternative
\begin{equation}
q_i(k,z)=\rho_i(z)q_i(k\mid z),
\label{eq:structured}
\end{equation}
in which each template the token considers carries its own Gaussian refinement of the state, and the token's belief becomes a genuine Gaussian mixture. It is an option recorded for later use; it is not implemented.

\paragraph{Motivation.}
In the evaluated projected predictive route, no step after the encoder writes the belief covariance: attention updates the mean and passes $\Sigma_i$ through, and the block MLP does the same. The covariance delivered to the block is then the token-type lookup $\Sigma(x_i)$, a function of the token identity alone, and it cannot register that ``bank'' is ambiguous in one context and resolved in another. The calibration probe saved by two recent $K=64$ runs on this route records exactly this: the covariance source is labeled input-token variance, and the across-token coefficient of variation of $\Tr\Sigma$ is $0.02$ to $0.03$. The structured posterior supplies a context-dependent covariance from the mean alone, because the mean does depend on context.

\paragraph{Evidence and prior.}
A structured posterior needs something to condition on. Treat the belief delivered by attention as a Gaussian message about the latent state, $\N(k;\mu_i,\Sigma_i)$, and the concept bank as the prior $\sum_a\pi_a\N(k;\nu_a,\Lambda_a)$. This is the reading used for conformity in Section~\ref{sec:social}, and it is the natural reading when attention outputs are messages rather than posteriors.

\begin{proposition}[Exact structured posterior]
\label{prop:structured}
With Gaussian evidence $\N(k;\mu_i,\Sigma_i)$ and the mixture prior of Eq.~\eqref{eq:mixture_prior}, the exact posterior over $(k,z)$ has the form of Eq.~\eqref{eq:structured} with
\begin{align}
\tilde\rho_{ia}&=\frac{\pi_a\N(\mu_i;\nu_a,\Lambda_a+\Sigma_i)}{\sum_b\pi_b\N(\mu_i;\nu_b,\Lambda_b+\Sigma_i)},
\label{eq:bayes_resp}\\
q_i(k\mid z=a)&=\N(k;m_{ia},C_{ia}),\qquad
C_{ia}=\bigl(\Lambda_a^{-1}+\Sigma_i^{-1}\bigr)^{-1},\qquad
m_{ia}=\mu_i+\Sigma_i(\Sigma_i+\Lambda_a)^{-1}(\nu_a-\mu_i).
\label{eq:component}
\end{align}
\end{proposition}

\begin{proof}
For each $a$, the product of two Gaussian densities in $k$ factorizes as $\N(k;\nu_a,\Lambda_a)\N(k;\mu_i,\Sigma_i)=\N(\mu_i;\nu_a,\Lambda_a+\Sigma_i)\N(k;m_{ia},C_{ia})$, with $C_{ia}$ and $m_{ia}$ as stated \citep{bishop2006}; the second form of $m_{ia}$ follows from $C_{ia}\Sigma_i^{-1}=\Lambda_a(\Lambda_a+\Sigma_i)^{-1}=I-\Sigma_i(\Sigma_i+\Lambda_a)^{-1}$. The joint density of $(k,z=a)$ given the evidence is proportional to $\pi_a$ times this product. Normalizing over $k$ and $a$ gives Eq.~\eqref{eq:bayes_resp} for the label marginal and Eq.~\eqref{eq:component} for the conditional.
\end{proof}

Each component mean $m_{ia}$ is the conformity update of Eq.~\eqref{eq:conformity}: the token's belief pulled toward template $a$ by the precision-weighted fraction $\Sigma_i(\Sigma_i+\Lambda_a)^{-1}$. The belief is therefore a mixture of ``what I would believe if template $a$ were the right reading,'' weighted by how well each reading predicts the evidence.

\paragraph{Two routing rules.}
Equation~\eqref{eq:bayes_resp} is not the rule of Proposition~\ref{prop:responsibilities}. Up to $a$-independent terms,
\begin{align*}
\log\tilde\rho_{ia}&=\log\pi_a-\tfrac12(\mu_i-\nu_a)^{\top}(\Lambda_a+\Sigma_i)^{-1}(\mu_i-\nu_a)-\tfrac12\log\det(\Lambda_a+\Sigma_i),\\
\log\rho^{*}_{ia}&=\log\pi_a-\tfrac12(\mu_i-\nu_a)^{\top}\Lambda_a^{-1}(\mu_i-\nu_a)-\tfrac12\Tr(\Lambda_a^{-1}\Sigma_i)-\tfrac12\log\det\Lambda_a.
\end{align*}
Both reduce to $\log\pi_a+\log\N(\mu_i;\nu_a,\Lambda_a)$ as $\Sigma_i\to0$, and they differ for uncertain beliefs. The mean-field rule treats $q_i$ as the posterior itself and asks whether template $a$ can cover the whole belief; the trace term penalizes sharp templates, so uncertain tokens route to broad ones. The Bayesian rule treats $q_i$ as noisy evidence and asks which template most plausibly generated the state; the evidence covariance is added to each template's, the distances shrink, and an uncertain token spreads its weight over the templates in proportion to their predictive density. Table~\ref{tab:structured} repeats the worked example of Section~\ref{sec:example} under Proposition~\ref{prop:structured}. The uncertain ``bank'' now splits $0.37$ and $0.37$ between the financial and riverside readings with $0.26$ on the generic template, where the mean-field rule gave $0.94$ to the generic template, and its component means $m_{i1}=(1.8,0.2)$ and $m_{i2}=(0.2,1.8)$ are the two senses. For polysemy, the Bayesian answer is arguably the better description: a use of ``bank'' is one of two senses, not a generic place. The choice between the rules is a modeling statement about what attention delivers, a posterior or a message.

\paragraph{Collapsing to one Gaussian.}
Attention, transport, and the gauge results all assume Gaussian beliefs. The structured posterior can be used inside the block and collapsed at its output by moment matching,
\begin{equation}
\mu_i^{+}=\sum_a\tilde\rho_{ia}m_{ia},
\qquad
\Sigma_i^{+}=\underbrace{\sum_a\tilde\rho_{ia}C_{ia}}_{\text{within}}
+\underbrace{\sum_a\tilde\rho_{ia}(m_{ia}-\mu_i^{+})(m_{ia}-\mu_i^{+})^{\top}}_{\text{between}},
\label{eq:collapse}
\end{equation}
which is the Gaussian closest to the mixture in $\KL(\text{mixture}\|\cdot)$. The between-template term is an ambiguity covariance: it is large exactly when the token splits its weight among templates whose refined means are far apart. It depends on context through $\mu_i$, so it gives the covariance a context-dependent update even on a route where nothing else writes it. In the toy example, $\Tr\Sigma^{+}$ has between-template part $0.96$ at ``bank'' and $0.02$ at confident ``deposit,'' a ratio of about $44$, while the incoming traces were $4$ and $0.6$.

\begin{table}[h]
\centering
\begin{tabular}{@{}lccccc@{}}
\toprule
Token state & $\tilde\rho_{1}$ & $\tilde\rho_{2}$ & $\tilde\rho_{3}$ & $\Tr(\text{within})$ & $\Tr(\text{between})$ \\
\midrule
``bank,'' $\mu=(1,1)$, $\Sigma=2I$ & $0.369$ & $0.369$ & $0.262$ & $1.22$ & $0.96$ \\
``deposit,'' $\mu=(1.8,0.2)$, $\Sigma=0.3I$ & $0.824$ & $0.015$ & $0.160$ & $0.40$ & $0.02$ \\
``deposit,'' $\mu=(1.8,0.2)$, $\Sigma=2I$ & $0.587$ & $0.163$ & $0.250$ & $1.20$ & $0.68$ \\
\bottomrule
\end{tabular}
\caption{The toy bank of Table~\ref{tab:toy_templates} under the exact structured posterior of Proposition~\ref{prop:structured}: Bayesian responsibilities of Eq.~\eqref{eq:bayes_resp} and the two parts of the collapsed covariance of Eq.~\eqref{eq:collapse}. Compare the mean-field responsibilities in Table~\ref{tab:toy_routing}.}
\label{tab:structured}
\end{table}

\paragraph{Values become exact.}
With joint templates $p_a(k,y)$ as in Eq.~\eqref{eq:gmr} and $y\mid k,a\sim\N(w_a+B_a(k-\nu_a),\cdot)$, the structured posterior gives the readout
\begin{equation}
\E[y\mid\text{evidence}]=\sum_a\tilde\rho_{ia}\bigl(w_a+B_a(m_{ia}-\nu_a)\bigr),
\label{eq:structured_readout}
\end{equation}
exactly, because the conditional mean is linear in $k$ within each component and the unobserved $y$ does not change the posterior over $k$. The mean-field readout of Section~\ref{sec:update} was exact only as $\Sigma_i\to0$; the structured version removes that restriction, and with $B_a=0$ it is the implemented readout with $\rho^{*}$ replaced by $\tilde\rho$.

\paragraph{Gauge covariance and cost.}
In reference coordinates the evidence density and each $\N(\mu_i;\nu_a,\Lambda_a+\Sigma_i)$ transform with the same Jacobian factor, which is independent of $a$ and cancels in Eq.~\eqref{eq:bayes_resp}, so the argument of Proposition~\ref{prop:covariant} carries over to the structured posterior. The cost depends on the covariance used. With diagonal $\Lambda_a$ and $\diag\Sigma_i$, every quantity in Eqs.~\eqref{eq:bayes_resp} and~\eqref{eq:component} is elementwise, $O(MK)$ per token as in the implemented block, and the diagonal of Eq.~\eqref{eq:collapse} is also $O(MK)$. A full between-template covariance costs $O(MK^2)$ per token through the low-rank form $D^{\top}\diag(\tilde\rho_i)D$ with $D$ the matrix of centered component means. With a full $\Sigma_i$, each template needs its own $(\Lambda_a+\Sigma_i)^{-1}$, which is $O(MK^3)$ per token and not practical at the widths used here.

\paragraph{Carrying mixtures through attention.}
The alternative to collapsing is to pass the mixture to the next layer. Attention would then compare mixtures, whose divergence has no closed form; the soft minimum of Proposition~\ref{prop:responsibilities} and its generalization by \citet{hershey2007} supply a variational surrogate, at $C^2$ component pairs per token pair for $C$-component beliefs, and the exact covariance statements of the attention operator would become statements about that surrogate. The collapsed version keeps every existing result and adds multimodality only where the model already has a mixture prior. It is the version to try first.

\paragraph{Diagnostics.}
The structured posterior makes ambiguity measurable per token and per context: the entropy of $\tilde\rho_i$, the between-template trace of Eq.~\eqref{eq:collapse}, and their correlation with next-token cross-entropy. These are the context-dependent counterparts of the calibration probe that the frozen covariance of the evaluated route cannot support.

\section{Templates as a shared cultural repertoire}
\label{sec:social}

The agent reading invites a social analogy: an individual reasoning within a conversation also consults a stock of shared expectations, norms, and scripts that no single conversation created. This section states what that analogy captures exactly, what it captures only in a modified variant, and where it fails. The short answer is that the templates behave like a shared cultural repertoire of norms and scripts that an individual consults, and not like a meta-agent in the sense of the multi-agent program.

\paragraph{What matches.}
Four features of the template bank have direct social counterparts. The bank is shared by every token in every sequence, as norms are shared by the members of a population. It changes on a slower timescale than the beliefs that consult it, as norms change across generations of cultural transmission rather than within a single exchange \citep{boyd1985}. It is consulted and not renegotiated: a token cannot alter a template during a forward pass, just as an individual does not rewrite a convention in the course of following it. And the quantities in the free energy of Eq.~\eqref{eq:free_energy} have natural readings. The divergence $\KL(q_i\|p_a)$ is the cost of an individual's belief deviating from norm $a$, the weight $\pi_a$ is the prevalence of the norm, and the width $\Lambda_a$ is its permissiveness: a broad template tolerates a wide range of beliefs, and a sharp one tolerates little.

\paragraph{Conformity and uncertainty.}
The tied-value construction of Section~\ref{sec:update} is the analog of conformity. Suppose agent $i$ fuses its belief with norm $a$ by precision weighting, treating $p_a$ as a Gaussian prior and its current belief as the evidence. The resulting conjugate Gaussian has mean
\begin{equation}
\mu_i^{+}=\mu_i+\Sigma_i(\Sigma_i+\Lambda_a)^{-1}(\nu_a-\mu_i),
\label{eq:conformity}
\end{equation}
of which the natural-gradient step Eq.~\eqref{eq:tied_update} is the first-order form when $\Sigma_i\ll\Lambda_a$. For isotropic $\Sigma_i=sI$ and $\Lambda_a=\lambda I$, the agent moves the fraction $s/(s+\lambda)$ of the way to the norm. Uncertain agents conform more, confident agents conform less, and strict norms pull harder than permissive ones. This is the pattern of informational social influence, in which others' judgments are taken as evidence \citep{deutsch1955}, and it is consistent with the classical observation that shared norms form readily when the stimulus is ambiguous and offers no external frame of reference \citep{sherif1936}. Combined with the routing result of the worked example, the mean-field model makes two predictions: an uncertain agent defers to broad, generic norms, and a confident one consults sharp, specific norms but moves little toward them. Under the Bayesian routing of Section~\ref{sec:structured} the first prediction changes: an uncertain agent hedges among the specific norms that could apply, in proportion to how well each explains its evidence, rather than retreating to a generic one. The implemented block uses untied values, so its update magnitude does not scale with $\Sigma_i$ in this way; only its routing does. Equation~\eqref{eq:conformity} describes the tied variant.

\paragraph{Norms versus scripts.}
Untied values correspond to a different social object. A value $w_a$ that is independent of the template mean does not say ``be more like the norm''; it says ``in a situation of kind $a$, this is what tends to come next.'' That is a script in the sense of \citet{schank1977}: a stored expectation about how a recognized kind of situation unfolds. A language model is trained to predict rather than to prescribe, so its trained templates should be read as scripts and schemas more than as norms of conduct. The two readings coexist in the joint-template construction of Eq.~\eqref{eq:gmr}, where $\nu_a$ says what the situation looks like and $w_a$ says what follows it.

\paragraph{How norms would form: the population average.}
The analogy suggests asking how a template would be set if the agents themselves determined it, rather than an external optimizer. Sum the free energy of Eq.~\eqref{eq:free_energy} over a population of beliefs $\{q_i\}$ with fixed responsibilities and minimize over the template. For a full-covariance template, the stationary point is
\begin{equation}
\pi_a=\frac1N\sum_i\rho^{*}_{ia},
\qquad
\nu_a=\frac{\sum_i\rho^{*}_{ia}\mu_i}{\sum_i\rho^{*}_{ia}},
\qquad
\Lambda_a=\frac{\sum_i\rho^{*}_{ia}\bigl[\Sigma_i+(\mu_i-\nu_a)(\mu_i-\nu_a)^{\top}\bigr]}{\sum_i\rho^{*}_{ia}},
\label{eq:mstep}
\end{equation}
the M-step of expectation--maximization for a Gaussian mixture fitted to beliefs rather than points \citep{dempster1977}. The diagonal case takes the diagonal of $\Lambda_a$. In social terms, a norm settles at the average belief of its adherents, weighted by how strongly each identifies with it; its prevalence is the fraction who adhere; and its permissiveness is the spread of their beliefs plus their own uncertainty, so a norm is strict exactly when its adherents are homogeneous and certain. This is bottom-up norm formation of the kind \citet{sherif1936} observed. Prevalence enters the routing only through $\log\pi_a$, so an agent's probability of adopting a norm, all else equal, is proportional to the norm's prevalence. That is unbiased transmission. Conformist transmission, in which common variants are adopted more than in proportion to their frequency \citep{henrich1998}, would correspond to replacing $\log\pi_a$ by $\kappa\log\pi_a$ with $\kappa>1$; the derivation fixes $\kappa=1$. The implemented model does not use Eq.~\eqref{eq:mstep}. It learns $(\nu_a,\Lambda_a,\pi_a,w_a)$ by gradient descent on next-token prediction, so trained templates are the norms that best serve prediction, which need not be the averages of their adherents.

\paragraph{Where the analogy fails.}
The templates are not meta-agents in the sense of the multi-agent program. There, a higher-level agent is formed from a particular coalition of constituent agents once they reach coherence, its state is built from theirs, and it acts back on them. A template is built from no particular tokens, exists before any sequence is read, and is not changed by the tokens that consult it during inference. The in-context counterpart of a meta-agent is a coherent cluster of tokens formed by attention, such as the tokens of one noun phrase agreeing on a referent; that structure lives in the peer channel, not in the template bank. The distinction is between a culture's repertoire, which templates resemble, and a particular group, which a meta-agent is. A second failure is that the model has informational influence but no normative influence. \citet{deutsch1955} separated conformity driven by evidence from conformity driven by the wish for approval, and the second requires sanctions that nothing in the model supplies; the only pressure toward a template is the prior cost. A third concerns the gate. Human beings hold several roles and norms at once, which fits the independent gates of Remark~\ref{rem:independent} better than the softmax, which asserts that one template explains the agent. The softmax reading is closer to identifying the single situation one is in than to the set of norms one obeys.

\section{Implementation}

The block is selected in \texttt{vfe3} by \texttt{block\_mlp\_mode="gaussian\_template"} with \texttt{use\_block\_mlp=True}; the default mode is unchanged. Each of the $L$ layers owns an untied bank of $M=\texttt{block\_mlp\_expansion}\times K$ concepts with parameters $\nu\in\mathbb{R}^{M\times K}$, $\log\lambda\in\mathbb{R}^{M\times K}$, $\log\pi\in\mathbb{R}^{M}$ (unnormalized; the softmax normalizes), and $w\in\mathbb{R}^{M\times K}$, for $3MK+M$ parameters per layer. At initialization $\lambda_{ak}=1$ and $\pi_a$ is uniform, so the templates share the covariance $I$ and Eq.~\eqref{eq:affine_logits} holds exactly for any belief covariance until training separates the widths; $\nu$ and $w$ are drawn with the same scales as the rows of $W_1$ and the columns of $W_2$ in the coordinate MLP. The logits are computed from Eq.~\eqref{eq:matmul_form}, whose first two terms are token-dependent contractions against $M\times K$ tables and whose third depends only on the parameters; the readout adds one more contraction. This is $3MK$ multiply-adds, or about $6MK$ FLOPs, per token per layer, against $2MK$ multiply-adds ($4MK$ FLOPs) for the MLP. Only \texttt{block\_mlp\_covariance="passthrough"} is supported. The settings \texttt{block\_mlp\_activation} and \texttt{block\_mlp\_covariance\_floor} are unread in this mode and are reported as inert when set to non-default values. At equal expansion the block has about $1.5$ times the parameters of the MLP, so a parameter-matched comparison uses about two thirds of the MLP's expansion.

The accompanying tests check the responsibilities against an exact full-covariance divergence, the shared-covariance reduction of Eq.~\eqref{eq:affine_logits}, agreement between full and diagonal covariance storage, finite-difference gradients with respect to the inputs and every parameter, nonzero finite gradients for every concept parameter inside a small model, and the reported parameter and FLOP counts.

\section{What this does and does not establish}

The derivation establishes that a block of the form of Eq.~\eqref{eq:implemented} follows from a mixture prior over learned Gaussian concepts under a mean-field variational posterior, that its routing weights are the same Gibbs form as attention with unit temperature, that it reduces to a softmax-gated MLP for templates with a shared covariance, and that its covariant form in reference coordinates has the same transformation law as the attention residual. It supplies an account of what the block computes, which the coordinate MLP lacks.

It does not establish that template routing matches or improves on the MLP in predictive loss. No training run has been performed. The nearest empirical precedent, persistent-memory attention, reported performance comparable to a feed-forward sublayer rather than better \citep{sukhbaatar2019}, and a softmax across units is a sparser, more competitive gate than independent activations, which may require more concepts to reach the same fit. The untied-value update is a declared variational readout, exact only in the limit of a concentrated belief. The implemented block is neither covariant in the sense of Proposition~\ref{prop:covariant} nor equipped with regression blocks, independent gates, or a covariance update; each is a stated extension. The structured posterior of Section~\ref{sec:structured} is exact under its message reading of attention, but it is not implemented, and its Bayesian routing differs from the mean-field routing of the implemented block for uncertain beliefs. The parameter-matched ablation of template routing against the GELU MLP, with variability across seeds, is the experiment that decides whether the derivation also buys performance. Useful diagnostics for that run are the entropy of $\rho^{*}_i$ across tokens and layers, the fraction of concepts that receive appreciable responsibility, and the correlation of routing with $\Tr\Sigma_i$, which tests the uncertainty-dependent behavior of the preceding section.

\begin{thebibliography}{99}

\bibitem[Bishop(2006)]{bishop2006}
C.~M.~Bishop.
\newblock \emph{Pattern Recognition and Machine Learning}.
\newblock Springer, New York, 2006.

\bibitem[Boyd and Richerson(1985)]{boyd1985}
R.~Boyd and P.~J.~Richerson.
\newblock \emph{Culture and the Evolutionary Process}.
\newblock University of Chicago Press, Chicago, 1985.

\bibitem[Dempster et~al.(1977)]{dempster1977}
A.~P.~Dempster, N.~M.~Laird, and D.~B.~Rubin.
\newblock Maximum likelihood from incomplete data via the {EM} algorithm.
\newblock \emph{Journal of the Royal Statistical Society, Series B}, 39(1):1--38, 1977.

\bibitem[Deutsch and Gerard(1955)]{deutsch1955}
M.~Deutsch and H.~B.~Gerard.
\newblock A study of normative and informational social influences upon individual judgment.
\newblock \emph{Journal of Abnormal and Social Psychology}, 51(3):629--636, 1955.

\bibitem[Elfwing et~al.(2018)]{elfwing2018}
S.~Elfwing, E.~Uchibe, and K.~Doya.
\newblock Sigmoid-weighted linear units for neural network function approximation in reinforcement learning.
\newblock \emph{Neural Networks}, 107:3--11, 2018.

\bibitem[Friston(2008)]{friston2008hierarchical}
K.~Friston.
\newblock Hierarchical models in the brain.
\newblock \emph{PLoS Computational Biology}, 4(11):e1000211, 2008.

\bibitem[Geva et~al.(2021)]{geva2021}
M.~Geva, R.~Schuster, J.~Berant, and O.~Levy.
\newblock Transformer feed-forward layers are key-value memories.
\newblock In \emph{Proceedings of the 2021 Conference on Empirical Methods in Natural Language Processing}, pages 5484--5495, 2021.

\bibitem[Ghahramani and Jordan(1994)]{ghahramani1994}
Z.~Ghahramani and M.~I.~Jordan.
\newblock Supervised learning from incomplete data via an {EM} approach.
\newblock In \emph{Advances in Neural Information Processing Systems 6}, pages 120--127, 1994.

\bibitem[Hendrycks and Gimpel(2016)]{hendrycks2016}
D.~Hendrycks and K.~Gimpel.
\newblock Gaussian error linear units ({GELUs}).
\newblock arXiv:1606.08415, 2016.

\bibitem[Henrich and Boyd(1998)]{henrich1998}
J.~Henrich and R.~Boyd.
\newblock The evolution of conformist transmission and the emergence of between-group differences.
\newblock \emph{Evolution and Human Behavior}, 19(4):215--241, 1998.

\bibitem[Hershey and Olsen(2007)]{hershey2007}
J.~R.~Hershey and P.~A.~Olsen.
\newblock Approximating the {K}ullback {L}eibler divergence between {G}aussian mixture models.
\newblock In \emph{IEEE International Conference on Acoustics, Speech and Signal Processing}, pages IV-317--IV-320, 2007.

\bibitem[Jacobs et~al.(1991)]{jacobs1991}
R.~A.~Jacobs, M.~I.~Jordan, S.~J.~Nowlan, and G.~E.~Hinton.
\newblock Adaptive mixtures of local experts.
\newblock \emph{Neural Computation}, 3(1):79--87, 1991.

\bibitem[Ramachandran et~al.(2017)]{ramachandran2017}
P.~Ramachandran, B.~Zoph, and Q.~V.~Le.
\newblock Searching for activation functions.
\newblock arXiv:1710.05941, 2017.

\bibitem[Ramsauer et~al.(2021)]{ramsauer2021}
H.~Ramsauer, B.~Sch{\"a}fl, J.~Lehner, P.~Seidl, M.~Widrich, T.~Adler, L.~Gruber, M.~Holzleitner, M.~Pavlovi{\'c}, G.~K.~Sandve, V.~Greiff, D.~Kreil, M.~Kopp, G.~Klambauer, J.~Brandstetter, and S.~Hochreiter.
\newblock Hopfield networks is all you need.
\newblock In \emph{International Conference on Learning Representations}, 2021.

\bibitem[Schank and Abelson(1977)]{schank1977}
R.~C.~Schank and R.~P.~Abelson.
\newblock \emph{Scripts, Plans, Goals and Understanding: An Inquiry into Human Knowledge Structures}.
\newblock Lawrence Erlbaum Associates, Hillsdale, NJ, 1977.

\bibitem[Sherif(1936)]{sherif1936}
M.~Sherif.
\newblock \emph{The Psychology of Social Norms}.
\newblock Harper \& Brothers, New York, 1936.

\bibitem[Sukhbaatar et~al.(2019)]{sukhbaatar2019}
S.~Sukhbaatar, E.~Grave, G.~Lample, H.~J{\'e}gou, and A.~Joulin.
\newblock Augmenting self-attention with persistent memory.
\newblock arXiv:1907.01470, 2019.

\end{thebibliography}

\end{document}
